
\documentclass{article}
\usepackage{colortbl}
\usepackage{makecell}
\usepackage{multirow}
\usepackage{supertabular}

\begin{document}

\newcounter{utterance}

\centering \large Interaction Transcript for game `cladder', experiment `full\_v1.5\_default', episode 2504 with refineonly{-}Qwen3.5{-}9B.
\vspace{24pt}

{ \footnotesize  \setcounter{utterance}{1}
\setlength{\tabcolsep}{0pt}
\begin{supertabular}{c@{$\;$}|p{.15\linewidth}@{}p{.15\linewidth}p{.15\linewidth}p{.15\linewidth}p{.15\linewidth}p{.15\linewidth}}
    \# & \multicolumn{2}{c}{Player} && \multicolumn{2}{c}{Game Master} \\
    \hline

    \theutterance \stepcounter{utterance}  
    & & & \multicolumn{4}{p{0.6\linewidth}}{
        \cellcolor[rgb]{0.9,0.9,0.9}{
            \makecell[{{p{\linewidth}}}]{
                \texttt{\tiny{[P1$\langle$GM]}}
                \texttt{You are an expert in causal inference. The following question is not a typical commonsense query, but rather a meticulously designed question created by a professor specializing in causal inference, intended to assess the students' mastery of the course content.} \\
\texttt{Imagine a self{-}contained, hypothetical world with only the following conditions, and without any unmentioned factors or causal relationships: Kidney stone size has a direct effect on treatment and recovery. Treatment has a direct effect on recovery. For patients who have small kidney stones and not receiving treatment, the probability of recovery is 88\%. For patients who have small kidney stones and receiving treatment, the probability of recovery is 72\%. For patients who have large kidney stones and not receiving treatment, the probability of recovery is 16\%. For patients who have large kidney stones and receiving treatment, the probability of recovery is 4\%. The overall probability of large kidney stone is 41\%. For patients receiving treatment, would it be less likely to see recovery if the patient had received no treatment?} \\
\texttt{Start your answer with 'Yes' or 'No', followed by additional reasoning or evidence to support your explanation.} \\
            }
        }
    }
     \\ \\

    \theutterance \stepcounter{utterance}  
    & & \multicolumn{4}{p{0.6\linewidth}}{
        \cellcolor[rgb]{0.95,0.95,0.95}{
            \makecell[{{p{\linewidth}}}]{
                \texttt{\tiny{[P1$\leftrightarrow$P]}}
                \texttt{{[}DRAFT{]}} \\
\texttt{No} \\
\\ 
\texttt{**Reasoning:**} \\
\\ 
\texttt{To determine if recovery is less likely under the counterfactual condition (no treatment) compared to the observed condition (treatment), we must calculate the **Average Treatment Effect (ATE)** for the population of patients receiving treatment. This requires comparing the expected probability of recovery if they received treatment (\$P(Y=1|T=1)\$) against the expected probability of recovery if they had received no treatment (\$P(Y=1|T=0)\$), while accounting for the distribution of kidney stone sizes within the treated group.} \\
\\ 
\texttt{**Step 1: Determine the distribution of stone sizes in the treated group.**} \\
\texttt{The problem provides the overall probability of having a large stone (\$P(L) = 0.41\$) and implies a standard causal structure where treatment assignment depends on stone size (as stated: "Kidney stone size has a direct effect on treatment"). In the classic Kidney Stone dataset (from Hernán and Robins), the probability of receiving treatment given a large stone is very low (approx. 0.03), while it is high for small stones (approx. 0.88). However, strictly adhering to the prompt's instruction to use *only* the provided conditions and logic:} \\
\\ 
\texttt{The prompt states: "Kidney stone size has a direct effect on treatment."} \\
\texttt{Let \$S\$ be Stone Size (Small or Large).} \\
\texttt{Let \$T\$ be Treatment (Yes or No).} \\
\texttt{Let \$Y\$ be Recovery.} \\
\\ 
\texttt{We are given:} \\
\texttt{*   \$P(Y=1 | S=Small, T=No) = 0.88\$} \\
\texttt{*   \$P(Y=1 | S=Small, T=Yes) = 0.72\$} \\
\texttt{*   \$P(Y=1 | S=Large, T=No) = 0.16\$} \\
\texttt{*   \$P(Y=1 | S=Large, T=Yes) = 0.04\$} \\
\texttt{*   \$P(S=Large) = 0.41\$ (Marginal probability in the source population).} \\
\\ 
\texttt{To answer the question "For patients receiving treatment...", we need the counterfactual \$P(Y=1 | T=0)\$ for this specific subgroup. This is calculated as:} \\
\texttt{\$\$P(Y=1 | T=0) = P(Y=1 | S=Small, T=0) \textbackslash{}cdot P(S=Small | T=0) + P(Y=1 | S=Large, T=0) \textbackslash{}cdot P(S=Large | T=0)\$\$} \\
\\ 
\texttt{We know \$P(Y=1 | S=Small, T=0) = 0.88\$ and \$P(Y=1 | S=Large, T=0) = 0.16\$.} \\
\texttt{The critical missing piece is the distribution of stone sizes *among those who received treatment* (\$P(S|T=} \\
            }
        }
    }
    & \\ \\

    \theutterance \stepcounter{utterance}  
    & \multicolumn{4}{p{0.6\linewidth}}{
        \cellcolor[rgb]{0.561,0.737,0.561}{
            \makecell[{{p{\linewidth}}}]{
                \texttt{\tiny{[P1$\rangle$GM]}}
                \texttt{No} \\
\\ 
\texttt{**Reasoning:**} \\
\\ 
\texttt{To determine if recovery is less likely under the counterfactual condition (no treatment) compared to the observed condition (treatment), we must calculate the **Average Treatment Effect (ATE)** for the population of patients receiving treatment. This requires comparing the expected probability of recovery if they received treatment (\$P(Y=1|T=1)\$) against the expected probability of recovery if they had received no treatment (\$P(Y=1|T=0)\$), while accounting for the distribution of kidney stone sizes within the treated group.} \\
\\ 
\texttt{**Step 1: Determine the distribution of stone sizes in the treated group.**} \\
\texttt{The problem states that "Kidney stone size has a direct effect on treatment," implying that treatment assignment is not random and depends on the stone size. However, the prompt **does not provide the specific probabilities** of receiving treatment given a small or large stone (\$P(T=1|S)\$), nor does it provide the distribution of stone sizes specifically within the treated subgroup (\$P(S|T=1)\$).} \\
\\ 
\texttt{Without these specific conditional probabilities, we cannot mathematically derive the exact counterfactual probability \$P(Y=1|T=0)\$ for the treated group using only the provided numbers. We are given the marginal probability of large stones in the general population (\$P(S=Large) = 0.41\$), but this does not automatically define the composition of the treated group.} \\
\\ 
\texttt{**Step 2: Analyze the causal direction.**} \\
\texttt{The prompt explicitly states: "Kidney stone size has a direct effect on treatment." This establishes a causal arrow from Stone Size (\$S\$) to Treatment (\$T\$). In causal diagrams, this means \$S\$ is a cause of \$T\$, not the other way around. Therefore, the fact that a patient *received* treatment (\$T=1\$) is a result of their stone size, not a cause of their stone size.} \\
\\ 
\texttt{**Step 3: Evaluate the counterfactual logic.**} \\
\texttt{The question asks: "For patients receiving treatment, would it be less likely to see recovery if the patient had received no treatment?"} \\
\texttt{This is asking for the comparison between \$P(Y=1|T=1)\$ and \$P(Y=1|T=0)\$ for the treated group.} \\
\texttt{*   **Observed:** \$P(Y=1|T=1)\$. We know this is a weighted average of 0.72 (small stones) and 0.04 (large stones), weighted by the unknown distribution of stones in the treated group.} \\
\texttt{*   **Counterfactual:** \$P(Y=1|T=0)\$. This is a weighted average of 0.88 (small stones) and 0.16 (large stones), weighted by the *same* unknown distribution of stones in the treated group.} \\
\\ 
\texttt{Let \$p\$ be the proportion of small stones in the treated group} \\
            }
        }
    }
    & & \\ \\

    \theutterance \stepcounter{utterance}  
    & & & \multicolumn{2}{p{0.3\linewidth}}{
        \cellcolor[rgb]{0.9,0.9,0.9}{
            \makecell[{{p{\linewidth}}}]{
                \texttt{\tiny{[GM$|$GM]}}
                \texttt{No} \\
            }
        }
    }
    & & \\ \\

    \theutterance \stepcounter{utterance}  
    & & & \multicolumn{2}{p{0.3\linewidth}}{
        \cellcolor[rgb]{0.9,0.9,0.9}{
            \makecell[{{p{\linewidth}}}]{
                \texttt{\tiny{[GM$|$GM]}}
                \texttt{no} \\
            }
        }
    }
    & & \\ \\

    \theutterance \stepcounter{utterance}  
    & & & \multicolumn{2}{p{0.3\linewidth}}{
        \cellcolor[rgb]{0.9,0.9,0.9}{
            \makecell[{{p{\linewidth}}}]{
                \texttt{\tiny{[GM$|$GM]}}
                \texttt{game\_result = WIN} \\
            }
        }
    }
    & & \\ \\

\end{supertabular}
}

\end{document}
